\section{Research Gaps}
\label{sec:gaps}

We state gaps as falsifiable absences in the literature.

\subsection{G1 --- No shared quality$\to$utility benchmark}

There is no community benchmark that (a)~publishes fixed random bit archives at multiple quality tiers, (b)~defines isolated stochastic-site hooks in reference trainers, and (c)~reports multi-seed metrics with pre-registered tests.
Existing papers each bring private streams and private tasks, preventing meta-analysis.

\subsection{G2 --- Missing isolation of stochastic sites}

Dropout-quality~\cite{koivu2022dropout} and initialisation-quality~\cite{bird2019softcomputing,heese2024biased} are studied separately; shuffle-quality and augmentation-quality almost not at all.
A full factorial (or at least one-site-at-a-time) design across S1--S4 (\S\ref{sec:background}) is absent.

\subsection{G3 --- Optimiser$\times$RNG interaction untested}

Despite rich theory on SGD versus Adam noise~\cite{zhou2020adamsgd,liu2020radam,wu2022alignment}, we found no controlled study that crosses optimiser family with measured bitstream quality of the shuffle/mask generator while holding all else fixed.
This is perhaps the highest-value gap for connecting the two literatures.

\subsection{G4 --- No task-relevant quality metric}

NIST/TestU01 failures do not map monotonically to accuracy deltas~\cite{antunes2025influence}.
Neural distinguishers~\cite{valle2024assessing} measure predictability, not optimiser utility.
Needed: metrics that predict downstream ML/optimiser sensitivity (e.g., estimated bias in mini-batch index distributions; autocorrelation at batch-scale lags; effective entropy consumed per epoch).

\subsection{G5 --- Scale mismatch}

Pilot studies use kilobytes to megabytes of random data.
Modern training can consume gigabytes of random bits across epochs (dropout masks alone).
Whether quality defects that appear only at Crush/BigCrush scales affect long training runs is unknown.
Large binary archives are uniquely suited to close this gap---yet remain under-used in published ML experiments.

\subsection{G6 --- Replication crisis on QRNG claims}

The Bird versus Heese disagreement~\cite{bird2019softcomputing,heese2024biased} is unresolved by direct, pre-registered replication on shared tasks with shared quality diagnostics.
Until then, QRNG-advantage claims should be treated as hypothesis-generating, not established.

\subsection{G7 --- Metaheuristics and deep learning are siloed}

PSO/GA/DE studies converged early on ``minimum quality suffices''~\cite{bastos2009pso,ma2014comparison}, while deep learning discussions rediscover similar ideas a decade later~\cite{antunes2025influence}.
Cross-domain synthesis---and tests of quasi-random versus pseudo-random \emph{mini-batch schedules}---are rare.

\subsection{G8 --- Reporting standards}

Papers seldom report: bytes consumed, wraparound policy, bit ordering, exact seeding maps, GPU determinism flags, or effect sizes with confidence intervals.
Without these, RNG-quality science cannot accumulate.

\subsection{Gap priority matrix}

\begin{table}[t]
\centering
\caption{Gap priorities for a bitstream-centred experimental programme.}
\label{tab:gap-priority}
\begin{tabular}{@{}clp{7.5cm}@{}}
\toprule
\textbf{ID} & \textbf{Priority} & \textbf{Rationale} \\
\midrule
G1 & High & Enables cumulative science \\
G2 & High & Removes attribution confounds \\
G3 & High & Links RNG quality to optimiser theory \\
G4 & Medium & Needs iterative metric design \\
G5 & High & Motivates large archives \\
G6 & Medium & Clarifies QRNG narrative \\
G7 & Medium & Broadens optimisation coverage \\
G8 & High & Cheap to adopt; high leverage \\
\bottomrule
\end{tabular}
\end{table}
